\chapter{Introduction}

\section{What Seal Is}

Seal is a fully end-to-end encrypted (E2EE) chat application. It offers direct
messages (DMs) and group channels, an attachment path for encrypted images, and
both a browser client and a native desktop client. The component documented in
this reference is the \emph{server}: a single, self-contained Rust binary named
\verb|seal-server| built on the \href{https://github.com/tokio-rs/axum}{Axum} web
framework, with persistence provided by \href{https://lancedb.github.io/lancedb/}{LanceDB},
an embedded columnar/vector database built on Apache Arrow.

The defining property of Seal is that \textbf{the server never sees plaintext}.
All encryption and decryption happens on the client using
\href{https://libsodium.gitbook.io/doc/}{libsodium}. The server stores and relays
opaque ciphertext, public keys (which are, by definition, public), and routing
metadata (who sent what to whom, and when). Private keys never leave the client.
This makes Seal a \emph{zero-knowledge relay}: a database or server compromise
exposes ciphertext and social graph metadata, but not message contents.

\section{Design Goals}

\begin{itemize}
  \item \textbf{Zero-knowledge server.} The server is a dumb, trustworthy pipe.
        It authenticates users, enforces channel membership, rate-limits abuse,
        and persists ciphertext --- nothing more. It holds no decryption keys.
  \item \textbf{Single static binary.} \verb|templates/|, \verb|static/|, and the
        default \verb|config.yaml| are compiled \emph{into} the binary via
        \verb|include_dir!| / \verb|include_str!|. The release artifact at
        \verb|target/release/seal-server| can be copied anywhere, given a
        \verb|JWT_SECRET| and a \verb|DATABASE_PATH|, and run.
  \item \textbf{Faithful port.} The server is a Rust port of an earlier FastAPI
        (Python) application. Many functions carry comments such as
        ``Mirrors \texttt{app/main.py}'' precisely because behaviour --- error
        strings, status codes, bcrypt hash format --- was kept wire-compatible
        with the Python original and its test suite.
  \item \textbf{Testability.} The bootstrap path is factored so integration tests
        spin up the \emph{real} router against a throwaway database on a random
        port. The suite exercises the HTTP and WebSocket surface end to end.
\end{itemize}

\section{Technology Stack}

\begin{center}
\begin{tabular}{@{}ll@{}}
\toprule
\textbf{Concern} & \textbf{Technology} \\
\midrule
HTTP / routing / WebSocket & Axum 0.8 on Tokio \\
Persistence & LanceDB 0.30 (Apache Arrow, \verb|arrow-array| 58) \\
Password hashing & bcrypt (\verb|DEFAULT_COST|) \\
Session tokens & JSON Web Tokens (HS256/384/512) \\
Client-side crypto & libsodium (X25519 + secretbox), in the browser/desktop \\
Config & YAML (\verb|serde_yml|) with environment-variable overrides \\
Asset bundling & \verb|include_dir| / \verb|include_str| \\
Logging & \verb|tracing| + \verb|tracing-subscriber| \\
Desktop client & Tauri v2 \\
\bottomrule
\end{tabular}
\end{center}

\section{A Note on Scope and Branch}

This reference documents the code on the \verb|main| branch. The repository also
carries a \verb|feat/object-storage| feature branch that teaches the database
layer to talk to S3, GCS, and Azure Blob Storage instead of the local
filesystem. That work is not yet merged into \verb|main|; on \verb|main| the
database connector is filesystem-only. Chapter~\ref{chap:object-storage}
describes the object-storage feature and clearly marks which parts live only on
the feature branch, so the picture is complete without misrepresenting what
\verb|main| actually compiles.

\section{How to Read This Book}

\begin{itemize}
  \item Chapter~\ref{chap:architecture} gives the 30{,}000-foot view: the
        three-layer model, \verb|AppState|, and a request's lifecycle.
  \item Chapters~\ref{chap:config}--\ref{chap:websocket} drill into the
        subsystems: configuration, the database, the REST API, and WebSockets.
  \item Chapter~\ref{chap:security} explains the cryptographic model and the
        server-side defences (auth, rate limiting, validation).
  \item Chapter~\ref{chap:source} is a module-by-module code walkthrough ---
        the reference you reach for when editing a specific file.
  \item Chapters~\ref{chap:testing}--\ref{chap:object-storage} cover the test
        suite, build and deployment, and the object-storage backend.
\end{itemize}

Throughout, inline identifiers such as \verb|db_ops::scan_where| and file paths
such as \verb|src/routes/messages.rs| point directly at the source so you can
jump from prose to code.
